% !PRINTED 131
\weberpage{131}
Wenn wir nun noch beweisen können, dass es in $U$ \gesperrt{genau $(p-1)\nu$ Elemente giebt, deren Grad durch $p$ theilbar ist,} so sind wir am Ziele; denn dann folgt, dass es unter den Elementen $U$ und also auch unter den Elementen von $P$ genau
\[
p\nu - (p-1)\nu = \nu
\]
giebt, deren Grad ein Theiler von $\nu$ ist.

Dies ergiebt sich aber aus Folgendem:

$\beta$) \gesperrt{Ist $v$ ein Element aus $V$, so gehören alle Elemente $vQ$ (deren Anzahl $p$ beträgt) zu $V$. Es giebt darunter $p-1$, deren Grad durch $p$ theilbar ist, und eines, dessen Grad nicht durch $p$ theilbar ist. In $V$ giebt es $(p-1)r_2$ Elemente, deren Grad durch $p$ theilbar ist, und $r_2$ Elemente, deren Grad nicht durch $p$ theilbar ist.}

Denn ist $a$ ein von $1$ verschiedenes Element von $Q$, so ist, da $v$ zu $R$ gehört, also $v^{-1}av$ in $Q$ enthalten ist, und da $Q$ aus den Potenzen von $a$ besteht,
\weq{1}{v^{-1}av = a^s.}
Durch wiederholte Anwendung ergiebt sich daraus für jeden Exponenten $h$
\weq{2}{v^{-h}av^h = a^{s^h}.}
Nehmen wir für $h$ den Grad des Elementes $v$, der nach der Voraussetzung ein Theiler von $p\nu$ ist, setzen also $v^h = 1$, so folgt
\weq{3}{a = a^{s^h}, \quad s^h \equiv 1 \ (\mathrm{mod}\ p).}
In $h$ gehen aber keine anderen Primfactoren auf als solche, die gleich oder grösser als $p$ sind, und also ist $p-1$ relativ prim zu $h$ und man kann der Congruenz
\weq{4}{xh \equiv 1 \ (\mathrm{mod}\ p-1)}
genügen; demnach ergiebt sich aus dem \gesperrt{Fermat'schen} Lehrsatze:
\weq{5}{s^{xh} \equiv s \equiv 1 \ (\mathrm{mod}\ p),}
also nach (1):
\weq{6}{av = va,}
d. h. $a$ und $v$ sind vertauschbar. Daraus folgt für jeden Exponenten $\lambda$
\weq{7}{(va)^\lambda = v^\lambda a^\lambda.}
